
%%%%%%%%%%%%%%%%%%%%%%%%%%%%%%%%%%%%%%%%
%           main body
%%%%%%%%%%%%%%%%%%%%%%%%%%%%%%%%%%%%%%%%

%-----------------------------------
%	TITLE PAGE
%-----------------------------------

\title[数学分析]{\texorpdfstring{\LARGE \bfseries 第九章\quad 数项级数}{第九章 数项级数}} % The short title appears at the bottom of every slide, the full title is only on the title page
\author[数项级数] % Your institution as it will appear on the bottom of every slide, maybe shorthand to save space
{\Large \bfseries 数项级数的收敛性}
% \author{Singular Value Decomposition} % Your name
% \date{\today} % Date, can be changed to a custom date
\date{}

%------------------------------------------------

\begin{document}

%------------------------------------------------

\begin{frame}

\titlepage % Print the title page as the first slide

\end{frame}

%------------------------------------------------

% \begin{frame}
% \frametitle{内容提要} % Table of contents slide, comment this block out to remove it
% \tableofcontents % Throughout your presentation, if you choose to use \section{} and \subsection{} commands, these will automatically be printed on this slide as an overview of your presentation
% \end{frame}

%-----------------------------------
%	PRESENTATION SLIDES
%-----------------------------------

%------------------------------------------------

\begin{frame}{芝诺悖论 (Zeno's Paradox)}

大约公元前425年，古希腊有一位名叫芝诺的学者，曾提出过若干个在数学发展是上产生过重大影响的悖论。“Achilles追赶乌龟”即是其中较著名的一个。

\vspace*{0.9em}

\begin{columns}
  \begin{column}{0.47\textwidth}
    \includegraphics[width=0.95\textwidth]{images/zeno-paradox.jpg}
  \end{column}
  \pause
  \begin{column}{0.51\textwidth}
    设乌龟在Achilles前面$S_1$米处向前爬行，Achilles在后面追赶。\\
    当Achilles用了$t_1$秒时间，跑完$S_1$米时，乌龟已向前爬了$S_2$米；\\
    当Achilles再用$t_2$秒时间，跑完$S_2$米时，乌龟又向前爬了$S_3$米；\\
    $\cdots\cdots$\\
    这样的过程可以一直继续下去，因此Achilles {\color{red}永远}也追不上乌龟。
  \end{column}
\end{columns}

\pause
\vspace*{0.9em}

% {\larger[2]\color{red}
% $$
% t_1 + t_2 + \cdots + t_n + \cdots = \text{无穷}?
% $$
% }
{\larger[2]
\begin{empheq}[box={\fancymath[colback=cyan!30,drop lifted shadow, sharp corners]}]{equation*}
  \color{red!80} t_1 + t_2 + \cdots + t_n + \cdots = \text{无穷}?
\end{empheq}
}

\end{frame}

%------------------------------------------------

\begin{frame}{一尺之捶，日取其半，万世不竭 |{\smaller[2] 《庄子·杂篇·天下》}}

\begin{center}

\includegraphics[width=0.85\textwidth]{images/一尺之捶-日取其半-万世不竭.jpg}
\end{center}

\pause
\vspace*{-0.9em}

% {\larger[2]\color{red}
% $$
% 1 = \dfrac{1}{2} + \dfrac{1}{4} + \cdots + \dfrac{1}{2^n} + \cdots ~ !
% $$
% }
{\larger[2]
\begin{empheq}[box={\fancymath[colback=cyan!30,drop lifted shadow, sharp corners]}]{equation*}
  \color{red!80} 1 = \dfrac{1}{2} + \dfrac{1}{4} + \cdots + \dfrac{1}{2^n} + \cdots ~ !
\end{empheq}
}

\pause

以上两个例子，我们遇到的都是无穷（{\color{red}可列}）多个数相加的情况，称之为{\color{red}“数项级数”}。那么数项级数有什么特点以及运算规律呢？这是我们接下来要学习的。

\end{frame}

%------------------------------------------------

\section[定义]{数项级数的定义}

%------------------------------------------------

\begin{frame}{数项级数定义}

设$x_1, x_2, \dots, x_n, \dots$是无穷可列个实数，我们称它们的{\color{red}“和”}
$$x_1 + x_2 + \cdots + x_n + \cdots$$
为\textbf{无穷数项级数}（简称\textbf{级数}），（形式上）记为$\sum\limits_{n=1}^\infty x_n,$ 其中$x_n$称为级数的\textbf{通项}或\textbf{一般项}。

\pause
\vspace*{0.5em}

\begin{block}{{\larger 问题} -- 如何进行无限求和？}
\vspace*{-1em}
{\larger[1]
\begin{align*}
{\color{red}\checkmark} & \phantom{+} x_1 + x_2 + \cdots + x_n \\
{\color{red}\textbf{?}} & \phantom{+} x_1 + x_2 + \cdots + x_n {\color{red} + \cdots}
\end{align*}
}
\end{block}

\end{frame}

%------------------------------------------------

\begin{frame}{部分和数列}

想法：我们考虑有限项（前$n$项）的和，然后将无限求和定义为其{\color{red}极限}！

\pause

\begin{align*}
  S_1 & = x_1, \\
  S_2 & = x_1 + x_2, \\
  S_3 & = x_1 + x_2 + x_3, \\
  \vdots & \\
  S_n & = x_1 + x_2 + \cdots + x_n, \\
  \vdots \\
  \rotatebox[origin=c]{-90}{\color{red}$\to$} \\
  & \hspace*{-1.5em} {\color{red}\sum\limits_{n=1}^\infty x_n}
\end{align*}

\end{frame}

%------------------------------------------------

\begin{frame}{作为部分和数列极限的无穷级数}

\begin{Def}
如果部分和数列$\{S_n\}$收敛于有限实数$S,$ 则称无穷级数$\sum\limits_{n=1}^\infty x_n$收敛，且称它的和为$S,$ 记为
$$S = \sum\limits_{n=1}^\infty x_n;$$
\pause
\vspace*{-0.3em}
如果部分和数列$\{S_n\}$发散，则称无穷级数$\sum\limits_{n=1}^\infty x_n$发散。
\end{Def}

\pause
\vspace*{0.3em}

由上述定义可知，只有当无穷级数收敛时，无穷多个实数的加法才是有意义的，并且它们的和就是级数的部分和数列的极限。所以，

\pause

\begin{empheq}[box={\fancymath[colback=cyan!30,drop lifted shadow, sharp corners]}]{equation*}
  \text{\bfseries\color{red!80} 级数的收敛与数列的收敛本质上是一回事！}
\end{empheq}

\end{frame}

%------------------------------------------------

\section[典型例子]{数项级数的典型例子}

%------------------------------------------------

\begin{frame}{数项级数的典型例子 | {\smaller 几何级数}}

设$|q| < 1,$ 则\textbf{几何级数}（即\textbf{等比数列}）
$$\sum\limits_{n=1}^\infty q^{n-1} = 1 + q + q^2 + \cdots + q^n + \cdots$$
是收敛的。
\pause
这是因为它的部分和数列通项为
$$S_n = \sum\limits_{k=1}^n q^{k-1} = \dfrac{1 - {\color{red}q^n}}{1 - q} ~~ \xrightarrow{\phantom{+} {\color{red}n \to \infty} \phantom{+}} ~~ \dfrac{1}{1 - q}$$

\end{frame}

%------------------------------------------------

\begin{frame}{数项级数的典型例子 | {\smaller 发散级数}}

级数
$$\sum\limits_{n=1}^\infty (-1)^{n-1} = 1 - 1 + 1 - \cdots + (-1)^{n-1} + \cdots$$
是发散的。
\pause
这是因为它的部分和数列通项为
$$S_n = \begin{cases}
0, & \text{$n$为偶数}, \\
1, & \text{$n$为奇数}.
\end{cases}$$
显然数列$\{S_n\}$是发散的。

\end{frame}

%------------------------------------------------

\begin{frame}{数项级数的典型例子 | {\smaller $\zeta$}}

$$\zeta(s) = \sum\limits_{n=1}^\infty \dfrac{1}{n^s} = 1 + \dfrac{1}{2^s} + \dfrac{1}{3^s} + \cdots + \dfrac{1}{n^s} + \cdots$$
被称作$\zeta$函数，是现代 \textbf{(解析) 数论}的中心研究对象之一。上述定义式对实部大于$1$的复数都有意义，并且可以解析延拓为整个复数域$\mathbb{C}$上的一个亚纯函数。著名的黎曼猜想就是关于$\zeta(s)$的：
% \begin{block}{黎曼猜想}
% \begin{center}
% $\zeta(s)$的所有非平凡零点的实部都等于$\dfrac{1}{2}.$
% \end{center}
% \end{block}
\begin{empheq}[box={\fancymath[colback=cyan!30,drop lifted shadow, sharp corners]}]{equation*}
  \text{\bfseries\color{red!80} $\zeta(s)$的所有非平凡零点的实部都等于$\dfrac{1}{2}.$}
\end{empheq}

\pause
\vspace*{-0.3em}

我们比较熟悉的是$s$是实数的情况。之前在学习数列极限的时候，我们已经证明了
\begin{center}
$s > 1$时， $\sum\limits_{n=1}^\infty \dfrac{1}{n^s}$收敛； \quad $s \leqslant 1$时， $\sum\limits_{n=1}^\infty \dfrac{1}{n^s}$发散到正无穷。
\end{center}

\pause
\vspace*{-0.3em}

这一级数在今后学习其他级数敛散性判别的时候也会有重要作用。

\end{frame}

%------------------------------------------------

\begin{frame}{芝诺悖论解惑}

设乌龟速度$v_1$ (米/秒) 与 Achilles 速度$v_2$之比为 $q ( 0 < q < 1).$\\
当 Achilles 跑完 $S_1$ 米时，乌龟已向前爬了$S_2 {\color{red} = qS_1}$米；\\
当 Achilles 跑完 $S_2$ 米时，乌龟又向前爬了$S_3 {\color{red} = qS_2 = q^2S_1}$米；\\
$\cdots\cdots$\\
当 Achilles 跑完 $S_n$ 米时，乌龟已向前爬了$S_{n+1} {\color{red} = q^nS_1}$米；\\
$\cdots\cdots$\\

\pause
\vspace*{0.5em}

将所有路程加起来，有
$$\color{red!80} S_1 + S_2 + \cdots + S_n + \cdots = S_1 + qS_1 + \cdots + q^{n-1}S_2 + \cdots = \dfrac{S_1}{1 - q}$$
这{\color{red}无限段}路程之和是{\color{red}有限}的,也就是说，当 Achilles 跑完路程$S = \dfrac{S_1}{1 - q}$ (即经过了{\color{red}有限}的时间$T = \dfrac{S_1}{(1-q)v_2}$秒) 时，他已经追上了乌龟。

\end{frame}

%------------------------------------------------

\section[基本性质]{数项级数的基本性质}

%------------------------------------------------

\begin{frame}{数项级数收敛的必要条件}

可以由数列的性质平行地导出数项级数的一些基本性质。

\pause

\begin{thm}[数项级数收敛的必要条件]
设级数$\sum\limits_{n=1}^\infty x_n$收敛，则其通项所构成的数列$\{x_n\}$是无穷小量，即
$$\lim\limits_{n\to\infty} x_n = 0.$$
\end{thm}

\pause

\textbf{证明：}
  设$\sum\limits_{n=1}^\infty x_n = S,$ 则对级数部分和 $S_n = \sum\limits_{k=1}^k x_k$ 成立
  $$\lim\limits_{n\to\infty} S_n = \lim\limits_{n\to\infty} S_{n-1} = S,$$
  于是得到
  $$\lim\limits_{n\to\infty} x_n = \lim\limits_{n\to\infty} (S_n - S_{n-1}) = \lim\limits_{n\to\infty} S_n - \lim\limits_{n\to\infty} S_{n-1} = 0.$$

\end{frame}

%------------------------------------------------

\begin{frame}{数项级数收敛的必要条件}

\begin{block}{注意}
$\lim\limits_{n\to\infty} x_n = 0$只是级数收敛的必要条件，并{\color{red}不是充分条件}。例如
$$\zeta(1) = \lim\limits_{n\to\infty} \dfrac{1}{n} = 1 + \dfrac{1}{2} + \dfrac{1}{3} + \cdots + \dfrac{1}{n} + \cdots$$
就满足这一必要条件，但它本身是发散的。
\end{block}

\pause
\vspace*{0.5em}

这一必要条件的一个重要作用在于，可以用于直接判断某些级数发散，例如我们已经看过的例子
$$\sum\limits_{n=1}^\infty (-1)^{n-1} = 1 - 1 + 1 - \cdots + (-1)^{n-1} + \cdots$$
现在可以直接通过其通项$(-1)^{n-1}$不是无穷小量得出这一级数发散的结论，而不需要再去求它的部分和数列，并判断其是否收敛。

\end{frame}

%------------------------------------------------

\begin{frame}{数项级数运算的线性性}

\begin{thm}[数项级数运算的线性性]
设$\sum\limits_{n=1}^\infty a_n = A,$ $\sum\limits_{n=1}^\infty b_n = B,$ $\alpha, \beta$是两个常数，那么
$$\sum\limits_{n=1}^\infty (\alpha a_n + \beta b_n) = \alpha A + \beta B.$$
\end{thm}

也就是说，对收敛级数可以进行加法以及数乘的运算。

\pause
\vspace{0.5em}

\textbf{证明：}
  设级数$\sum\limits_{n=1}^\infty a_n, \sum\limits_{n=1}^\infty b_n$的部分和数列分别为$\{S_n^{(1)}\}, \{S_n^{(2)}\},$ 那么对$\sum\limits_{n=1}^\infty (\alpha a_n + \beta b_n)$的部分和数列$\{S_n\}$有$S_n = \alpha S_n^{(1)} + \beta S_n^{(2)}.$ 对$n$取极限，并利用极限的线性性，即有
  $$\lim\limits_{n\to\infty} S_n = \alpha \lim\limits_{n\to\infty} S_n^{(1)} + \beta \lim\limits_{n\to\infty} S_n^{(2)} = \alpha A + \beta B.$$

\end{frame}

%------------------------------------------------

\begin{frame}{数项级数运算线性性的应用}

\begin{eg}
\begin{center}
求级数$\sum\limits_{n=1}^\infty \dfrac{4^{n+1} - 3 \cdot 2^n}{5^n}$的值。
\end{center}
\end{eg}

\pause
\vspace*{0.5em}

因为几何级数$\sum\limits_{n=0}^\infty \left( \dfrac{4}{5} \right)^n$与$\sum\limits_{n=0}^\infty \left( \dfrac{2}{5} \right)^n$都收敛，所以依据数项级数运算的线性性，有
\begin{align*}
  \sum\limits_{n=1}^\infty \dfrac{4^{n+1} - 3 \cdot 2^n}{5^n} & = \dfrac{16}{5} \sum\limits_{n=0}^\infty \left( \dfrac{4}{5} \right)^n - \dfrac{6}{5} \sum\limits_{n=0}^\infty \left( \dfrac{2}{5} \right)^n \\
  & = \dfrac{16}{5} \cdot \dfrac{1}{1 - \frac{4}{5}} - \dfrac{6}{5} \cdot \dfrac{1}{1 - \frac{2}{5}} \\
  & = 14.
\end{align*}

\end{frame}

%------------------------------------------------

\begin{frame}{数项级数运算的加法结合律}

待写

\end{frame}

%------------------------------------------------

\section{总结}

%------------------------------------------------

\begin{frame}{总结}

{\larger[2] \color{red} \ding{43}} 待写

\pause
\vspace{0.8em}

{\larger[2] \color{red} \ding{43}} 待写

\pause
\vspace{0.8em}

{\larger[2] \color{red} \ding{43}} 待写

\end{frame}

%------------------------------------------------

%------------------------------------------------
% % closing page
% \begin{frame}
% \HUGE{\centerline{\bfseries {\color{gray} The} {\color{zkblue} End}}}

% \pause

% \vspace{1.2em}

% \Large{\centerline{\bfseries {\color{gray}Thank} \quad {\color{zkblue} You}}}

% \end{frame}

%----------------------------------------------------------------------------------------

\end{document}
